\subsection{The exact arrow minimum at five objects}

\begin{theorem}[Proved]\label{A7:model32}
There is a finite acyclic category with five objects and thirty-two
arrows, identities included, whose classifying space is homotopy
equivalent to $\mathbb{CP}^2$.
\end{theorem}
\begin{proof}
Here is its complete composition law, with no auxiliary monoid or
unlisted choices. The objects are $0<1<2<3<4$ and their identities
are arrows $0,1,2,3,4$, respectively. The nonidentity hom-sets are
\[
\begin{array}{c|l@{\qquad}c|l}
01&5,6&02&7,8,9\\
03&10,11&04&12\\
12&13,14,15,16&13&17,18,19,20\\
14&21,22&23&23,24,25,26\\
24&27,28,29&34&30,31.
\end{array}
\]
There are no other arrows. An identity has its usual composites.
Every composite with source $0$ and target $4$ is arrow $12$.
For each remaining object triple $i<j<k$, the matrix $T_{ijk}$
below gives the arrow $b\circ a$, with rows $a\in C(i,j)$ and
columns $b\in C(j,k)$ in the increasing orders just displayed:
\[
\begin{aligned}
T_{012}&=\begin{pmatrix}7&7&8&8\\8&9&7&9\end{pmatrix},&
T_{013}&=\begin{pmatrix}10&10&11&10\\10&11&10&10\end{pmatrix},\\[3pt]
T_{023}&=\begin{pmatrix}10&10&10&11\\10&10&11&10\\10&11&10&10\end{pmatrix},&
T_{123}&=\begin{pmatrix}
17&17&18&19\\17&18&20&19\\20&20&19&18\\17&18&19&17
\end{pmatrix},\\[3pt]
T_{124}&=\begin{pmatrix}21&21&22\\21&22&22\\22&22&21\\21&22&21\end{pmatrix},&
T_{134}&=\begin{pmatrix}21&21\\21&22\\22&21\\22&22\end{pmatrix},\\[3pt]
T_{234}&=\begin{pmatrix}27&27\\27&28\\28&29\\29&27\end{pmatrix}.&&
\end{aligned}
\]
Together with the singleton-target rule these are all ten strict
object-triple composition tables. Every displayed entry is typed.
For associativity, triples involving an identity are automatic.
The strict object quadruples $0124,0134,0234$ have target $04$,
so both bracketings equal $12$. The remaining comparisons are
exactly
\[
 T_{023}(T_{012}(a,b),c)
   =T_{013}(a,T_{123}(b,c)),\qquad
 T_{124}(b,T_{234}(c,d))
   =T_{134}(T_{123}(b,c),d).
\]
Here the arguments denote the actual arrow labels in the tables.
Direct substitution verifies each equality in all $32$ cases.
The rules therefore define a category. Its support and identity
endomorphisms make it acyclic. Its hom profile is
\[
 (2,3,2,1,4,4,2,4,3,2),
\]
whose sum is $27$; adding the five identities gives $32$ arrows.

We supply integral cycles and a cocycle. On chronological
two-simplices $(a,b)$, let $\alpha$ be zero except as follows;
when a row lists several arrows on each side it means all of
the indicated composable pairs:
\[
\begin{array}{c|l|r}
a&b&\alpha(a,b)\\\hline
5&15,16,17,19,21&1\\
7,15&23,26,27&1\\
8,13&24,28&-1\\
11,18,24&30&-1
\end{array}
\]
Write
\[
 (r_{13},r_{14},r_{15},r_{16})=(1,-2,-1,2),\qquad
 (s_{23},s_{24},s_{25},s_{26})=(3,-1,-1,-1).
\]
Define
\[
\begin{split}
y&=\sum_{b=13}^{16}r_b\bigl([5,b]-[6,b]\bigr),\\
z&=\sum_{b=13}^{16}\sum_{c=23}^{26}r_b s_c
 \bigl([5,b,c,30]-[5,b,c,31]
       -[6,b,c,30]+[6,b,c,31]\bigr).
\end{split}
\]
These data satisfy
\[
 \delta\alpha=0,\qquad \partial y=0,\qquad \partial z=0,
 \qquad \langle\alpha,y\rangle
       =\langle\alpha\smile\alpha,z\rangle=1.
\]
For the cycle identities, $\sum_b r_b=\sum_c s_c=0$.
The $r$-weighted composites in $T_{012}$ agree for rows $5$ and
$6$: both give $-[7]+[8]$. The $s$-weighted composites in
$T_{234}$ agree for columns $30$ and $31$: both give
$2[27]-[28]-[29]$. Finally
$\sum_{b,c}r_bs_c[T_{123}(b,c)]=0$. These three equalities,
together with the two outer difference factors in $z$, verify
all five four-boundary components. Substitution in the
three-boundary formula verifies $\delta\alpha=0$ on all $128$
three-simplices. The two pairings are also immediate from the
display: the first is $r_{15}+r_{16}=1$, and the second is this
same value times $-s_{24}=1$.

The complete normalized integral nerve has cell counts
\[
 (f_0,f_1,f_2,f_3,f_4)=(5,27,89,128,64).
\]
Its boundary ranks are $(4,23,65,63)$, and every nonzero Smith
invariant of all four boundaries is $1$. These are finite exact
integer calculations on the fully specified category, reproduced
independently by
\path{acyclic/code/certify_literal32_integral.py} and
\path{monoid/code/check_acyclic_literal_integral.py}.
The former constructs chains by strictly increasing object
sequences and finds integral witnesses by clearing and normalizing
rational-kernel columns, accepting only directly verified unit
pairings. The latter constructs chains recursively and uses
unimodular Hermite transformations to obtain saturated integral
kernels and a primitive cohomology class. Neither imports the
search or the other checker. Both independently verify all $481$
typed associativity equations, including identities, and the
boundary-square identities. Their complete matrices and witnesses
are supplied with
\path{acyclic/outputs/cp2_five_32_integral_certificate.json} and
\path{monoid/certificates/acyclic32_candidate0003_independent.verified.json}.
The former is also checked against the seven displayed tables and
the displayed $\alpha,y,z$ by the separate transcription checker
\path{acyclic/code/check_model32_display.py}.

As explained in Theorem~\ref{A7:model47}, the Smith invariants
of the full boundaries determine torsion in homology because
the kernel of a homomorphism between free abelian groups is
saturated. Thus the integral homology is
\[
 H_*(BC;\mathbb Z)=(\mathbb Z,0,\mathbb Z,0,\mathbb Z),
\]
and is zero above degree four. The unit pairings above make
$[\alpha]$ and $[\alpha]^2$ integral generators.

For simple connectivity, choose an arrow $p_i:0\to i$ and
$q_i:i\to4$ for every $i$, choosing identities at the endpoints.
For every arrow $f:i\to j$, the singleton $C(0,4)$ gives
$q_jfp_i=q_jp_j$. In the groupoid completion this implies
$f=p_jp_i^{-1}$, so every loop is trivial. Equivalently, these
are consequences of the literal edge-and-triangle presentation
of the nerve. Both integral checkers verify the required
composite identities for every arrow. The simply connected CW
classification used in Theorem~\ref{A7:model47}, namely
\cite[Proposition~4C.3, pp.~430--431]{CP2-Hatcher2002}, now gives
$BC\simeq\mathbb{CP}^2$.
\end{proof}

\begin{corollary}[Proved]\label{A7:exact-object-and-arrow-minima}
The least object count among finite acyclic models of
$\mathbb{CP}^2$ is $5$, and the least arrow count among those
five-object models, including identities, is exactly $32$.
In the notation of Part A.7,
\[
 \boxed{\nu(\mathbb{CP}^2)=5,\qquad a_5(\mathbb{CP}^2)=32.}
\]
\end{corollary}
\begin{proof}
The dimension bound of Proposition~\ref{A7:dimension-chain}
requires five objects. Theorem~\ref{A7:through31} excludes every
five-object acyclic category with at most thirty-one arrows,
and Theorem~\ref{A7:model32} attains both remaining bounds.
The complete thirty-two-arrow census is not needed for this
minimality conclusion: its upper witness is one literal category,
and the lower bound is the separately completed through-thirty-one
exclusion.
\end{proof}
